Un protó es mou en direcció positiva de l'eix $OY$ en una regió on existeix un camp elèctric $\vec{E} = 3,0\times 10^{5}\,\vec{k}\ \mathrm{N/C}$ i un camp magnètic $\vec{B} = 0,60\,\vec{\imath}\ \mathrm{T}$.

\begin{apartat}{1,25}
Feu una representació esquemàtica de les forces que actuen sobre el protó indicant clarament els eixos, direccions i sentits. En quines condicions el protó no es desvia? Justifiqueu la resposta.
\end{apartat}

\begin{apartat}{1,25}
Un electró que es mou amb una velocitat $\vec{v} = 5\times 10^{5}\,\vec{\jmath}\ \mathrm{m/s}$ entra en aquesta regió. L'electró es desviarà? En cas afirmatiu, indiqueu cap a quina direcció es desvia. Justifiqueu la resposta representant esquemàticament les forces que actuen sobre l'electró.
\end{apartat}
